\section{Matrices}

\SourcePage{48}{51}
Assume for convenience that the system in question possesses a discrete
energy spectrum (every relation derived below extends directly to the
continuous-spectrum case as well). Let
$\Psi=\sum_n a_n\Psi_n$ denote the expansion of an arbitrary wave function over
the stationary-state wave functions $\Psi_n$. If this expansion is
substituted into the definition (3.8) of the mean value of some quantity $f$, we
obtain
\begin{equation}
  \bar f=\sum_n\sum_m a_n^*a_m f_{nm}(t),
  \tag{11.1}
\end{equation}
where $f_{nm}(t)$ denote the integrals
\begin{equation}
  f_{nm}(t)=\int\Psi_n^*\hat f\Psi_m\,dq.
  \tag{11.2}
\end{equation}
The set of quantities $f_{nm}(t)$ for all possible $n$, $m$ is called the
\textit{matrix} of the quantity $f$, and each $f_{nm}(t)$ is spoken of as a
\textit{matrix element}, corresponding to the transition from state $m$ to
state $n$\footnote{The matrix representation of physical quantities was introduced
by \textit{Heisenberg} (\textit{W. Heisenberg}) in 1925, even before
Schr\"{o}dinger's discovery of the wave equation. ``Matrix mechanics''
was subsequently developed by \textit{Born, Heisenberg and Jordan}
(\textit{M. Born, P. Jordan}).}.

The time dependence of the matrix elements $f_{nm}(t)$ is determined (if the
operator $\hat f$ does not contain $t$ explicitly) by the time dependence of the
functions $\Psi_n$. Substituting for them the expressions (10.1),
\SourcePage{49}{52}
we find that
\begin{equation}
  f_{nm}(t)=f_{nm}e^{i\omega_{nm}t},
  \tag{11.3}
\end{equation}
where
\begin{equation}
  \omega_{nm}=\frac{E_n-E_m}{\hbar}
  \tag{11.4}
\end{equation}
is the transition frequency between states $n$ and $m$, and the quantities
\begin{equation}
  f_{nm}=\int\psi_n^*\hat f\psi_m\,dq
  \tag{11.5}
\end{equation}
constitute the time-independent matrix of the quantity $f$, which is what one
ordinarily uses\footnote{In connection with the indeterminacy of the phase
factor in the normalized wave functions (see \secref{2}), the matrix elements
$f_{nm}$ (and $f_{nm}(t)$) are likewise defined only to within factors of the
form $\exp[i(\alpha_m-\alpha_n)]$. Here too this indeterminacy does not affect
the physical results.}.

One obtains the matrix elements of the derivative $\dot f$ by differentiating
those of the quantity $f$ with respect to time; this is an immediate consequence
of the relation
\begin{equation}
  \bar{\dot f}=\dot{\bar f}
  =\sum_n\sum_m a_n^*a_m\dot f_{nm}(t).
  \tag{11.6}
\end{equation}
In view of (11.3) we thus have for the matrix elements of $\dot f$:
\begin{equation}
  \dot f_{nm}(t)=i\omega_{nm}f_{nm}(t)
  \tag{11.7}
\end{equation}
or (cancelling the time factor $\exp(i\omega_{nm}t)$ on both sides) for the
time-independent matrix elements:
\begin{equation}
  (\dot f)_{nm}=i\omega_{nm}f_{nm}
  =\frac{i}{\hbar}(E_n-E_m)f_{nm}.
  \tag{11.8}
\end{equation}
For simplicity of notation in the formulas, we derive below all relations for
the time-independent matrix elements; exactly the same relations hold also for
matrices that depend on time.

For the matrix elements of the complex conjugate quantity $f^*$ of $f$, taking
into account the definition of the adjoint operator, we obtain
\[
  (f^*)_{nm}=\int\psi_n^*\hat f^+\psi_m\,dq
  =\int\psi_n^*\tilde{f^*}\psi_m\,dq
  =\int\psi_m\hat f^*\psi_n^*\,dq,
\]
i.e.
\begin{equation}
  (f^*)_{nm}=(f_{mn})^*.
  \tag{11.9}
\end{equation}
\SourcePage{50}{53}
For real physical quantities, which are usually all that we
consider, we therefore have
\begin{equation}
  f_{nm}=f_{mn}^*.
  \tag{11.10}
\end{equation}
Such matrices, like the operators corresponding to them, are called
\textit{Hermitian}.

Matrix elements with $n=m$ are called \textit{diagonal}. These elements do not
depend on time at all, and from (11.10) it is clear that they are real. The
element $f_{nn}$ is the mean value of $f$ in the state $\Psi_n$.

The rule of \textit{matrix multiplication} is readily derived. As a preliminary
step, observe that the formula
\begin{equation}
  \hat f\psi_n=\sum_m f_{mn}\psi_m.
  \tag{11.11}
\end{equation}
holds. This is nothing other than the expansion of the function $\hat f\psi_n$
in the functions $\psi_m$ with coefficients determined according to the general
rule (3.5). Bearing this formula in mind, we write
\[
  \hat f\hat g\psi_n=\hat f(\hat g\psi_n)
  =\sum_k g_{kn}\hat f\psi_k=\sum_{k,m}g_{kn}f_{mk}\psi_m.
\]
Since, on the other hand, we must have
$\hat f\hat g\psi_n=\sum_m(fg)_{mn}\psi_m$, we obtain
\begin{equation}
  (fg)_{mn}=\sum_k f_{mk}g_{kn}.
  \tag{11.12}
\end{equation}
This rule matches the standard mathematical prescription for multiplying
matrices: one multiplies the rows of the first factor by the columns of
the second.

Specifying a matrix is equivalent to specifying the operator itself. In
particular, it makes it possible in principle to determine the eigenvalues of a
given physical quantity and the eigenfunctions corresponding to them.

Let us consider the values of all quantities at some definite instant of time
and expand an arbitrary wave function $\Psi$ in the eigenfunctions of the
Hamiltonian, i.e. in the time-independent wave functions $\psi_m$ of the
stationary states:
\begin{equation}
  \Psi=\sum_m c_m\psi_m,
  \tag{11.13}
\end{equation}
\SourcePage{51}{54}
where the expansion coefficients are denoted by $c_m$. Let us substitute this
expansion into the equation $\hat f\Psi=f\Psi$, which determines the
eigenvalues and eigenfunctions of the quantity $f$. We have
\[
  \sum_m c_m\hat f\psi_m=f\sum_m c_m\psi_m.
\]
Multiply this equation on both sides by $\psi_n^*$ and integrate over $dq$.
Each of the integrals $\int\psi_n^*\hat f\psi_m\,dq$ on the left-hand side of
the equality is the corresponding matrix element $f_{nm}$. On the right-hand
side, all integrals $\int\psi_n^*\psi_m\,dq$ with $m\ne n$ vanish by virtue of
the orthogonality of the functions $\psi_m$, and
$\int\psi_n^*\psi_n\,dq=1$ by virtue of their
normalization\footnote{In accordance with the general rule (\secref{5}), the set of
coefficients $c_n$ of the expansion (11.13) can be regarded as the wave function
in the ``energy representation'' (the variable here being the index $n$,
which enumerates the energy eigenvalues). The matrix $f_{nm}$ then plays the
role of the operator $\hat f$ in this representation, whose action on the wave
function is given by the expression on the left-hand side of equation (11.14).
The formula $\bar f=\sum_n\sum_m c_n^*(f_{nm}c_m)$ then corresponds to the
general formula for the mean value of a quantity expressed through its operator
and the wave function of the state in question.}:
\begin{equation}
  \sum_m f_{nm}c_m=fc_n,
  \tag{11.14}
\end{equation}
or $\sum_m(f_{nm}-f\delta_{nm})c_m=0$, where
$\delta_{nm}=1$ for $n=m$ and $\delta_{nm}=0$ for $n\ne m$.

Thus we have obtained a system of homogeneous linear algebraic equations (with
unknowns $c_m$). As is well known, such a system has nonzero solutions only
if the determinant formed from the coefficients in the equations vanishes, i.e.
under the condition
\begin{equation}
  |f_{nm}-f\delta_{nm}|=0.
  \tag{11.15}
\end{equation}
The roots of this equation (where $f$ plays the role of the unknown) give
the possible values of the quantity $f$. The set of quantities $c_m$ satisfying
equations (11.14) upon setting $f$ equal to one of these values yields the
corresponding eigenfunction.
\SourcePage{52}{55}
Suppose that in the definition (11.5) for the matrix elements of $f$ we adopt,
as the functions $\psi_n$, the eigenfunctions of this very quantity; owing to the
equation $\hat f\psi_n=f_n\psi_n$ we obtain
\[
 f_{nm}=\int\psi_n^*\hat f\psi_m\,dq
 =f_m\int\psi_n^*\psi_m\,dq.
\]
Since the functions $\psi_m$ are orthogonal and normalized,
it follows that $f_{nm}=0$ for $n\ne m$ and $f_{mm}=f_m$. Thus only diagonal matrix
elements turn out to be nonzero, each of them being equal to the corresponding
eigenvalue of the quantity $f$; a matrix in which only these elements are
nonzero is said to be reduced to \textit{diagonal form}. In particular, in the
ordinary representation---where the stationary-state wave functions serve
as the $\psi_n$---the energy matrix is diagonal (along with the matrices of
all other physical quantities that have definite values in the stationary
states). More broadly, when the matrix of a quantity $f$ is constructed using
the eigenfunctions of an operator $g$, one speaks of it as the matrix of $f$ in the
representation in which $g$ is diagonal. Everywhere where it is not otherwise
specified, by the matrix of a physical quantity we shall henceforth understand
the matrix in the ordinary representation, in which the energy is diagonal.
All that has been said above concerning the time dependence of matrices refers,
of course, only to this ordinary representation\footnote{Keeping in mind that
the energy matrix is diagonal, one readily verifies that the equality (11.8) is
the operator relation (9.2) written in matrix form.}.

The matrix representation of operators makes it possible to prove the
theorem cited in \secref{4}: if two operators commute, they possess
a common complete system of eigenfunctions. Let $\hat f$ and $\hat g$ be two
such operators. From $\hat f\hat g=\hat g\hat f$ and the rule of matrix
multiplication (11.12) it follows that
\[
  \sum_k f_{mk}g_{kn}=\sum_k g_{mk}f_{kn}.
\]
Taking as the system of functions $\psi_n$, in terms of which the matrix
elements are computed, the eigenfunctions of the operator $\hat f$, we have
$f_{mk}=0$ for $m\ne k$, so that the written equality reduces to
$f_mg_{mn}=g_{mn}f_n$, or $g_{mn}(f_m-f_n)=0$. If all the eigenvalues
$f_n$ of the quantity $f$ are distinct, then for all $m\ne n$ we have
$f_m-f_n\ne0$, so that we must have $g_{mn}=0$.
\SourcePage{53}{56}
Thus the matrix $g_{mn}$ likewise turns out to be diagonal, i.e. the functions
$\psi_n$ are eigenfunctions also of the physical quantity $g$. If, however,
among the values $f_n$ there are equal ones, then the matrix elements $g_{mn}$
corresponding to each such group of functions $\psi_n$ will, in general, be
nonzero. However, linear combinations of the functions $\psi_n$ corresponding
to one eigenvalue of the quantity $f$ are likewise its eigenfunctions; one can
always choose these combinations in such a way as to make the corresponding
nondiagonal matrix elements $g_{mn}$ vanish, and thus in this case too we
obtain a system of functions that are eigenfunctions simultaneously for the
operators $\hat f$ and $\hat g$.

Let us note a useful formula
\begin{equation}
  \left(\frac{\partial H}{\partial\lambda}\right)_{nn}
  =\frac{\partial E_n}{\partial\lambda},
  \tag{11.16}
\end{equation}
where $\lambda$~--- a parameter on which the Hamiltonian $\hat H$ depends
(and with it the energy eigenvalues $E_n$). Indeed, differentiating the
equation $(\hat H-E_n)\psi_n=0$ with respect to $\lambda$ and then multiplying
it from the left by $\psi_n^*$, we obtain
\[
  \psi_n^*(\hat H-E_n)\frac{\partial\psi_n}{\partial\lambda}
  =\psi_n^*\left(\frac{\partial E_n}{\partial\lambda}
  -\frac{\partial\hat H}{\partial\lambda}\right)\psi_n.
\]
On integrating over $dq$, the left-hand side of this equality vanishes, since
\[
 \int\psi_n^*(\hat H-E_n)\frac{\partial\psi_n}{\partial\lambda}\,dq
 =\int\frac{\partial\psi_n}{\partial\lambda}(\hat H-E_n)^*\psi_n^*\,dq
\]
by virtue of the Hermiticity of the operator $\hat H$. The right-hand side
gives the desired equality.

In the modern literature, a system of notation (introduced by
\textit{Dirac}) is widely used, in which matrix elements are written as\footnote{We
will use both ways of denoting matrix elements in this book. The notation
(11.17) is especially convenient when each index has to be written as a
combination of several letters.}
\begin{equation}
  \langle n|f|m\rangle.
  \tag{11.17}
\end{equation}
This symbol is composed of the notation for the quantity $f$ and the symbols
$|m\rangle$ and
\SourcePage{54}{57}
$|n\rangle$, denoting respectively the initial and final states. If
there is a complete set of state functions $|n_1\rangle,|n_2\rangle,\ldots$, then
the expansion coefficients of the wave function of the state $|m\rangle$ are
denoted by $\langle n_i|m\rangle$:
\begin{equation}
  \langle n_i|m\rangle=\int\psi_{n_i}^*\psi_m\,dq.
  \tag{11.18}
\end{equation}